% !PRINTED 362
\weberpage{362}

If therefore one puts
\weq{5}{a_1 = N(\mathfrak{m}_1) = \mathfrak{m}_1\mathfrak{m}_1',\quad \frac{b_1+\sqrt{D}}{2} \equiv 0\ (\mathrm{mod}\ \mathfrak{m}_1),}
then one obtains another quadratic form $(a_1,b_1,c_1)$, of which we can show that it is equivalent to $(a,b,c)$.

We put, for brevity:
\[
\omega = \frac{b+\sqrt{D}}{2a},\quad \omega_1 = \frac{b_1+\sqrt{D}}{2a_1}
\]
and remark that
\[
\eta a_1 = \eta \mathfrak{m}_1\mathfrak{m}_1' = \mathfrak{m}\mathfrak{m}_1',
\]
\[
\eta \omega_1 a_1 = \frac{\mathfrak{m}}{\mathfrak{m}_1}\,\frac{b_1+\sqrt{D}}{2}
\]
are numbers in $[Q]$ divisible by $\mathfrak{m}$ [by (5)]. Accordingly, by \S\,97, 7., the rational integers $\alpha$, $\beta$, $\gamma$, $\delta$ can be so determined that:
\weq{6}{\begin{aligned}\eta a_1&=a(\alpha-\gamma\omega),\\ \eta a_1\omega_1&=a(-\beta+\delta\omega),\end{aligned}}
and one proves exactly as in \S\,95 that
\[
\alpha\delta - \beta\gamma = 1
\]
must hold, and hence that the two forms
\[
(a,b,c),\quad (a_1,b_1,c_1)
\]
are equivalent. We have therefore:

11. \gesperrt{Theorem. Every ideal class with respect to $[Q]$ that is represented by an ideal $\mathfrak{a}$ corresponds to a class of forms $A$ represented by the form $(a,b,c)$, and conversely.}

In order to obtain the class of forms $A$, one takes a primary ideal $\mathfrak{m}$ of the class $A^{-1}$ reciprocal to $A$ and puts $a = N(\mathfrak{m})$,
$\frac{b+\sqrt{D}}{2} \equiv 0\ (\mathrm{mod}\ \mathfrak{m})$. And if conversely the form $(a,b,c)$ is given, then the greatest common divisor $\mathfrak{m}$ of
\[
a,\quad \frac{b+\sqrt{D}}{2}
\]
is an ideal of the class $A^{-1}$.

\begin{center}
\S\,100. \gesperrt{Ideal classes with respect to the orders.}
\end{center}

The concept of the equivalence of ideals, and the theorems on class numbers, can be presented most clearly if
%CONT>>
